% paper/sections/08_0_bf_failure.tex
%
% §8.2.0  BF 失敗モード 5 種（SP-J §1）
%         §8.2（Balanced-Force）冒頭に配置し、
%         以降の P-1..P-7 原則（§8.3）の動機付けとする。
%
% ═══════════════════════════════════════════════════════════════════════════════
\subsection{Balanced-Force 失敗モード 5 種（SP-J §1）}
\label{sec:bf_failure_modes}
% ═══════════════════════════════════════════════════════════════════════════════

連続問題での静止液滴平衡
\begin{equation}
  -\bnabla p + \mathbf{f}_\sigma = \mathbf{0}
  \quad\text{（Young--Laplace: }[p]_\Gamma=\sigma\kappa\text{）}
  \label{eq:bf_continuous_balance}
\end{equation}
を\textbf{離散的に}保持するには
\begin{equation}
  -G_h p + f_{\sigma,h} \approx \mathbf{0}
  \label{eq:bf_discrete_balance}
\end{equation}
が要請される．微分階数ではなく\textbf{離散演算子の整合性}が
寄生流れ（spurious currents）を支配する．
SP-J \cite{SPJ} §1 は，実装で頻出する 5 つの破綻モードを列挙する：

\begin{tcolorbox}[warnbox,title={F-1..F-5：Balanced-Force 失敗モード}]
\begin{description}[leftmargin=2em,style=nextline]
  \item[\textbf{F-1} 評価位置不一致]
    $G_h p$ と $f_{\sigma,h}$ が異なる格子位置（節点 vs 面）で評価される．
    $(G_h p)_f \neq (f_\sigma)_f$ が定義上保証できない．
  \item[\textbf{F-2} PPE 勾配と補正勾配の不一致]
    PPE 内部で用いる $G_h$ と，$\mathbf{u}^* \to \mathbf{u}^{n+1}$ 補正で用いる $G_h$ が異なる．
    ある方程式を解き，別の方程式で補正していることになる．
  \item[\textbf{F-3} $D_h$ と $G_h$ が離散共役でない]
    $D_h \neq -G_h^\ast$．PPE 演算子が SPD でなくなり，
    離散エネルギー恒等式を崩し，CG が適用できない．
  \item[\textbf{F-4} $\betaf=(1/\rho)_f$ の経路不一致]
    PPE / 速度補正 / 界面張力の 3 経路で $\betaf$ の平均化
    （算術 vs 調和）が異なる．高密度比 $\rho_l/\rho_g\gg 1$ で
    補正残差が界面上で $\Ord{1}$ となる．
  \item[\textbf{F-5} 界面ジャンプ無補正の CCD ステンシル]
    $[p]_\Gamma=\sigma\kappa$ を跨ぐ CCD ステンシルが
    滑らかな場として扱われ，ジャンプが数値的に平滑化される．
    GFM/IIM による面フラックス補正が必要．
\end{description}
\end{tcolorbox}

\noindent\textbf{重要な事実}：
F-1..F-5 は\textbf{すべて微分階数と独立}である．
「高次 CCD を全量に適用」した素朴実装は F-1..F-5 のどれか（多くの場合複数）を
必ず踏む．CHK-172 の rising bubble PoC では F-2 と F-4 が同時発生し，
PPE RHS の発散だけを修正しても補正側が別 $G_h$ を参照する限り
BF 残差は減らない（SP-J §2.P-2 結論）．

\noindent\textbf{方針}：
§\ref{sec:bf_seven_principles} で 7 つの設計原則 P-1..P-7 を定式化し，
§\ref{sec:fccd_bf_sub_system} で FCCD 面ジェットによる
$\Ord{H^4}$ の BF 残差抑制を示す．

% ═══════════════════════════════════════════════════════════════════════════════
